% Lösungen Einheit 5 von 6 – Transformationen (zusammenbau v0.1)
\einheitenkopf{Einheit 5 von 6 · Transformationen}

\erg{19}{a) innen: x-Richtung \quad b) Verschiebung um $3$ nach oben \quad c) Verschiebung um $2$ nach links \quad d) Scheitel $(1 | -2)$, Form wie $f$ \quad e) $g(x) = (x - 3)^3 - 12(x - 3) + 5$; $H'(1 | 21)$ \quad f) $\frac{7{,}5}{3} = 2{,}5$ \quad g) richtig: $f'(2) = g'(2) = 12$; die Verschiebung nach oben ändert keine Steigung \quad h) Amplitude $= 3$, Periode $= \pi$, Mittellinie $y = 1$ \quad i) $g(x) = 0{,}1x^2 - 2$; Höhe in der Mitte $f(0) - g(0) = 4$ dm \quad j) $g'(5) = f'(3) = 4$ \quad k) $W_h = ]-10; 5]$: die Verschiebung in x-Richtung ändert die Wertemenge nicht; Spiegelung und Streckung mit $3$ machen aus $[-1; 4[$ das Intervall $]-12; 3]$, die Verschiebung um $2$ nach oben ergibt $]-10; 5]$}

19d)

\begin{ksys}[xmin=-4,xmax=4,ymin=-3,ymax=5,klein]\parabel{1}{0}{0}{f}\parabel{1}{1}{-2}{g}\end{ksys}

\erg{20}{a) Fehler: $+3$ im Argument wirkt gegenläufig. Richtig: Verschiebung um $3$ nach links \quad b) $g$ nimmt an der Stelle $x + 2$ den Wert an, den $f$ an der Stelle $x$ hat; jeder Punkt wandert also um zwei nach rechts. \quad c) um $2$ nach rechts und $1$ nach unten: $g(x) = (x - 2)^2 - 1$ \quad d) Periode $= 4\pi \approx 12{,}57$ h: so lange dauert es von Hochwasser zu Hochwasser; höchster Wasserstand $= 7$ m, niedrigster $= 3$ m}
